\documentclass[12pt, twoside]{article}
\input{../preamble}

\fancyhead[LE]{\thepage}
\fancyhead[RO]{Name: \hspace{4cm} \,\\ \hspace{2.9cm} \,}
\fancyhead[LO]{La Scuola d'Italia / Dr.\ Huson / IB Math: Kinematics \\* 11 March 2026}

\begin{document}
\subsubsection*{5.3 Classwork: Interpreting Motion Graphs}

\begin{enumerate}[itemsep=0.4cm]

%--- Problem 1: Displacement-time graph ---
\item The displacement $s$ (metres) of a particle from the origin is shown
      for $0 \leq t \leq 10$.

\begin{center}
\vspace{-0.3cm}
\begin{tikzpicture}[scale=0.5]
  \draw[-{Stealth}] (-0.5,0) -- (11,0) node[right] {$t$ (s)};
  \draw[-{Stealth}] (0,-4.5) -- (0,6.5) node[above] {$s$ (m)};
  \foreach \x in {1,...,10} \draw[gray!30] (\x,-4.5)--(\x,6.5);
  \foreach \y in {-4,...,6} \draw[gray!30] (-0.5,\y)--(11,\y);
  \foreach \x in {2,4,6,8,10} \draw (\x,0.1)--(\x,-0.1) node[below]{\small\x};
  \foreach \y in {-4,-2,2,4,6} \draw (0.1,\y)--(-0.1,\y) node[left]{\small\y};
  \node[below left] at (0,0) {\small 0};
  \draw[very thick]
    (0,2) -- (3,6) -- (6,6) -- (8,-2) -- (10,-4);
  \foreach \pt in {(0,2),(3,6),(6,6),(8,-2),(10,-4)}
    \fill \pt circle (3pt);
\end{tikzpicture}
\vspace{-0.3cm}
\end{center}

\begin{enumerate}[itemsep=0.1cm]
  \item Write down the initial displacement of the particle. \hfill [1]
  \item Find the velocity during $0 \leq t \leq 3$. \hfill [2]
  \item Describe the motion of the particle during $3 \leq t \leq 6$. \hfill [1]
  \item Find the velocity during $6 \leq t \leq 8$. \hfill [2]
  \item At what time does the particle pass through the origin? \hfill [2]
  \item Find the total \textbf{distance} travelled from $t = 0$ to $t = 10$. \hfill [3]
  \item Find the \textbf{displacement} from $t = 0$ to $t = 10$. \hfill [1]
\end{enumerate}


%--- Problem 2: Matching graphs ---
\item The graphs below show the displacement and velocity of a particle.
      Match each description (i)--(iv) to the correct time interval. \hfill [4]

\begin{center}
\begin{tikzpicture}[scale=0.5]
  \draw[-{Stealth}] (-0.5,0) -- (9,0) node[right] {$t$};
  \draw[-{Stealth}] (0,-1.5) -- (0,5.5) node[left] {$s$};
  \foreach \x in {2,4,6,8} \draw (\x,0.1)--(\x,-0.1) node[below]{\small\x};
  \draw[very thick]
    (0,0) -- (2,4) -- (4,4) -- (6,2) -- (8,2);
  \foreach \pt in {(0,0),(2,4),(4,4),(6,2),(8,2)}
    \fill \pt circle (3pt);
  \node[above] at (4.5,5.5) {\textbf{Displacement--time graph}};
\end{tikzpicture}
\hspace{1.5cm}
\begin{tikzpicture}[scale=0.5]
  \draw[-{Stealth}] (-0.5,0) -- (9,0) node[right] {$t$};
  \draw[-{Stealth}] (0,-2.5) -- (0,4.5) node[above] {$v$};
  \foreach \x in {2,4,6,8} \draw (\x,0.1)--(\x,-0.1) node[below]{\small\x};
  \draw[very thick]
    (0,2) -- (2,2) -- (2,0) -- (4,0) -- (4,-1) -- (6,-1) -- (6,0) -- (8,0);
  \node[above] at (4.5,4.5) {\textbf{Velocity--time graph}};
\end{tikzpicture}
\end{center}

\begin{enumerate}[label=(\roman*), itemsep=0.1cm]
  \item The particle is moving forward at constant speed.
  \item The particle is stationary.
  \item The particle is moving backward at constant speed.
  \item The particle is stationary (for the second time).
\end{enumerate}

\newpage

%--- Problem 3: Velocity-time graph ---
\item The velocity $v$ (m\,s$^{-1}$) of a cyclist is shown for $0 \leq t \leq 12$ seconds.

\begin{center}
\vspace{-0.3cm}
\begin{tikzpicture}[scale=0.55]
  \draw[-{Stealth}] (-0.5,0) -- (13,0) node[right] {$t$ (s)};
  \draw[-{Stealth}] (0,-3.5) -- (0,7.5) node[above] {$v$ (m\,s$^{-1}$)};
  \foreach \x in {1,...,12} \draw[gray!30] (\x,-3.5)--(\x,7.5);
  \foreach \y in {-3,...,7} \draw[gray!30] (-0.5,\y)--(13,\y);
  \foreach \x in {2,4,6,8,10,12} \draw (\x,0.1)--(\x,-0.1) node[below]{\small\x};
  \foreach \y in {-2,2,4,6} \draw (0.1,\y)--(-0.1,\y) node[left]{\small\y};
  \node[below left] at (0,0) {\small 0};
  \draw[very thick]
    (0,0) -- (4,6) -- (8,6) -- (10,0) -- (12,-2);
  \foreach \pt in {(0,0),(4,6),(8,6),(10,0),(12,-2)}
    \fill \pt circle (3pt);
\end{tikzpicture}
\vspace{-0.3cm}
\end{center}

\begin{enumerate}[itemsep=0.5cm]
  \item Find the acceleration during $0 \leq t \leq 4$. \hfill [2]
  \item Describe the motion during $4 \leq t \leq 8$. \hfill [1]
  \item Find the acceleration during $8 \leq t \leq 10$. \hfill [2]
  \item Find the total distance travelled from $t = 0$ to $t = 10$.

  \emph{Hint: the distance is the area under the velocity curve.} \hfill [3]
  \vspace{1.5cm}
  \item The cyclist continues with constant acceleration after $t = 10$.
        What does the negative velocity during $10 < t \leq 12$ tell you
        about the cyclist's motion? \hfill [1]
  \item Find the total distance travelled from $t = 0$ to $t = 12$. \hfill [2]
  \item Find the average speed of the cyclist over the 12-second interval. \hfill [2]
\end{enumerate}

\end{enumerate}
\end{document}
